\chapter*{KATA PENGANTAR}
\addcontentsline{toc}{chapter}{KATA PENGANTAR}

Puji dan syukur penulis panjatkan kepada Tuhan Yang Maha Esa atas berkat dan
rahmat-Nya sehingga laporan tugas akhir yang berjudul \textit{"Simulasi
	Operasional Armada Bus Listrik untuk Perencanaan Infrastruktur SPKLU pada
	Elektrifikasi Penuh TransJogja"} dapat diselesaikan sebagai salah satu syarat
kelulusan tingkat sarjana. Penulis menyadari bahwa penyelesaian tugas akhir
ini tidak terlepas dari dukungan banyak pihak. Oleh karena itu, penulis ingin
mengucapkan terima kasih kepada:

\begin{enumerate}
	\item Bapak Ir. I Gusti Bagus Baskara Nugraha, S.T., M.T., Ph.D. selaku
	dosen pembimbing dan pengampu mata kuliah II4092 atas segala bentuk
	masukan, kritik, dan saran serta waktu yang telah diberikan selama
	pengerjaan tugas akhir.
	
	\item Ibu Dr. Fetty Fitriyanti Lubis, S.T., M.T. selaku dosen mata kuliah
	II4092 atas arahannya dalam pelaksanaan tugas akhir di Program Studi
	Sistem dan Teknologi Informasi ITB.
	
	\item Bapak Ahmad Izzuddin, S.T., M.T. selaku Dosen Penguji Sidang
	Proposal yang telah memberikan masukan dan saran.
	
	\item Seluruh dosen pengampu mata kuliah yang telah memberikan ilmu
	pengetahuan dan pengalaman selama penulis berkuliah di Institut
	Teknologi Bandung.
	
	\item Bapak, Ibu, dan keluarga penulis yang selalu memberikan doa,
	dukungan, dan motivasi selama berkuliah hingga menyelesaikan tugas
	akhir.
	
	\item Wanda Dantini Putri, pacar penulis, yang senantiasa hadir
	dengan kesabaran, perhatian, dan dukungan tanpa henti, serta menjadi
	salah satu sumber semangat terbesar bagi penulis dalam menghadapi
	setiap tantangan selama menyelesaikan Tugas Akhir ini.
	
	\item Dzikri, Fadhlan, dan Iskandar, teman satu kelompok tugas akhir
	penulis, yang telah berproses bersama dan saling mendukung hingga
	Tugas Akhir ini dapat diselesaikan.
	
	\item Teman-teman lain di Institut Teknologi Bandung yang telah
	memberikan penulis pengalaman berharga selama menjalani perkuliahan.
\end{enumerate}

\begin{flushright}
	Bandung, 26 Agustus 2026\\
	
	Penulis \\
	\vspace{1cm}
	Muhammad Hanif Al Faithoni
\end{flushright}